\section{Introduction à l'IA générative}

\begin{frame}{Qu'est-ce que l'IA générative ?}

\textbf{Définition}

\vspace{0.3cm}

L'IA générative est une branche de l'intelligence artificielle qui \textbf{crée de nouveaux contenus} (textes, images, sons) à partir de données existantes.

\vspace{0.3cm}

\onslide<2->
\textbf{Différence clé :}
\begin{itemize}
    \item Systèmes traditionnels : analysent ou classifient
    \item \textcolor{blue}{Systèmes génératifs : produisent du contenu original}
\end{itemize}

\vspace{0.3cm}

\onslide<3>
\textbf{Exemples :}
\begin{itemize}
    \item \textbf{Génération de texte :} articles, chatbots
    \item \textbf{Génération d'images :} œuvres d'art (DALL-E, Stable Diffusion)
\end{itemize}

\end{frame}

\begin{frame}{Multimodalité}

\textbf{L'IA générative peut traiter différents types de données :}

\vspace{0.3cm}

\begin{itemize}
    \item<+-> \textbf{Texte $\rightarrow$ Image}
        \begin{itemize}
            \item Créer une image à partir d'une description textuelle
        \end{itemize}

    \vspace{0.2cm}
    \item<+-> \textbf{Image $\rightarrow$ Texte}
        \begin{itemize}
            \item Description automatique d'images
            \item Exemple : "une image montre un chien jouant dans un parc"
        \end{itemize}

    \vspace{0.2cm}
    \item<+-> \textbf{Autres modalités}
        \begin{itemize}
            \item Génération de voix ou de musique à partir de texte
        \end{itemize}
\end{itemize}

\vspace{0.3cm}

\onslide<4>
\begin{exampleblock}{Exemple concret}
Prompt : \textit{"un paysage montagneux au coucher de soleil, style impressionniste"} $\rightarrow$ génération d'une image correspondante
\end{exampleblock}

\end{frame}

\begin{frame}{Capacités}

\textbf{Génération de texte :}
\begin{itemize}
    \item Texte cohérent et contextuel
    \item Rédaction, traduction, résumé
\end{itemize}

\vspace{0.3cm}

\onslide<2->
\textbf{Génération d'images :}
\begin{itemize}
    \item Images haute résolution
    \item Respect de prompts complexes
\end{itemize}

\vspace{0.3cm}

\onslide<3>
\textbf{Applications variées :}
\begin{itemize}
    \item Aide à la rédaction
    \item Création artistique
    \item Assistance médicale
    \item Et bien d'autres...
\end{itemize}

\end{frame}

\begin{frame}{Technologies sous-jacentes}

\begin{enumerate}
    \item<+-> \textbf{Transformers}
        \begin{itemize}
            \item Architecture basée sur l'attention
            \item Capture le contexte dans les séquences
            \item Traite texte et images
        \end{itemize}

    \vspace{0.2cm}
    \item<+-> \textbf{Modèles de diffusion}
        \begin{itemize}
            \item Technique récente
            \item Images de haute qualité
            \item Exemple : Stable Diffusion
        \end{itemize}

    \vspace{0.2cm}
    \item<+-> \textbf{Grandeur des données}
        \begin{itemize}
            \item Énormes quantités de données nécessaires
            \item Entraînement intensif
        \end{itemize}
\end{enumerate}

\end{frame}

\begin{frame}{Techniques d'entraînement}

\begin{itemize}
    \item<+-> \textbf{Fine-tuning}
        \begin{itemize}
            \item Adapter un modèle pré-entraîné à une tâche spécifique
            \item Exemple : entraîner sur des données médicales
        \end{itemize}

    \vspace{0.2cm}
    \item<+-> \textbf{Apprentissage par renforcement}
        \begin{itemize}
            \item Améliorer via récompenses
            \item Feedback humain pour améliorer la qualité
        \end{itemize}

    \vspace{0.2cm}
    \item<+-> \textbf{Modèles multimodaux}
        \begin{itemize}
            \item Intégrer plusieurs types de données
            \item Texte + image pour tâches complexes
        \end{itemize}
\end{itemize}

\end{frame}


\begin{frame}{Bénéfices de l'IA générative}

\begin{itemize}
    \item<+->[\colormark{green}] \textbf{Productivité accrue}
        \begin{itemize}
            \item[\textcolor{white}{x}] Automatisation des tâches répétitives ou créatives
        \end{itemize}

    \vspace{0.2cm}
    \item<+->[\colormark{green}] \textbf{Accessibilité}
        \begin{itemize}
            \item[\textcolor{white}{x}] Démocratisation de la création de contenu
            \item[\textcolor{white}{x}] Pas de compétences techniques requises
        \end{itemize}

    \vspace{0.2cm}
    \item<+->[\colormark{green}] \textbf{Applications innovantes}
        \begin{itemize}
            \item[\textcolor{white}{x}] Aide à la recherche scientifique
            \item[\textcolor{white}{x}] Éducation personnalisée
            \item[\textcolor{white}{x}] Nouvelles formes de création
        \end{itemize}
\end{itemize}

\end{frame}

\begin{frame}{Risques de l'IA générative}

\begin{itemize}
    \item<+->[\colorwarning{red}] \textbf{Désinformation}
        \begin{itemize}
            \item[\textcolor{white}{x}] Génération de faux contenus
            \item[\textcolor{white}{x}] Deepfakes, articles fake news
        \end{itemize}

    \vspace{0.2cm}
    \item<+->[\colorwarning{red}] \textbf{Biais et discrimination}
        \begin{itemize}
            \item[\textcolor{white}{x}] Reproduction des biais des données d'entraînement
            \item[\textcolor{white}{x}] Exemple : stéréotypes de genre
        \end{itemize}

    \vspace{0.2cm}
    \item<+->[\colorwarning{red}] \textbf{Propriété intellectuelle}
        \begin{itemize}
            \item[\textcolor{white}{x}] Qui détient les droits sur un contenu généré par IA ?
        \end{itemize}

    \vspace{0.2cm}
    \item<+->[\colorwarning{red}] \textbf{Impact sur l'emploi}
        \begin{itemize}
            \item[\textcolor{white}{x}] Automatisation de tâches humaines
        \end{itemize}
\end{itemize}

\end{frame}

\begin{frame}{Qu'est-ce que l'AGI ?}
\framesubtitle{Artificial General Intelligence}

\textbf{Quid de l'AGI :}

\vspace{0.3cm}

Système d'IA capable de \textbf{raisonner et d'apprendre dans n'importe quel domaine}, avec une intelligence comparable à celle d'un humain.

\vspace{0.5cm}

\onslide<2>
\textbf{Différence avec l'IA générative :}

\begin{itemize}
    \item \textcolor{red}{IA générative :} spécialisée (texte, images)
    \item \textcolor{blue}{AGI :} intelligence générale, polyvalente
\end{itemize}

\end{frame}

\begin{frame}{État actuel de la recherche sur l'AGI}

\begin{itemize}
    \item<+->\colorwarning{red} \textbf{Pas de consensus scientifique}
        \begin{itemize}
            \item Faisabilité : débattue
            \item Proximité : opinions divergentes
        \end{itemize}

    \vspace{0.2cm}
    \item<+->\colorwarning{red} \textbf{Influences économiques}
        \begin{itemize}
            \item Débats influencés par investissements
            \item Course technologique
            \item Considérations économiques $>$ preuves scientifiques
        \end{itemize}

    \vspace{0.2cm}
    \item<+->\colorwarning{red} \textbf{Opinions des experts divisées}
        \begin{itemize}
            \item Certains : encore très loin
            \item D'autres : progrès rapides possibles
        \end{itemize}
\end{itemize}

\end{frame}

\begin{frame}{Message clé sur l'AGI}

\begin{block}{Points importants}
\begin{itemize}
    \item L'IA générative est \textcolor{red}{loin d'atteindre l'AGI}
    \item \textcolor{red}{Pas de consensus scientifique} sur :
        \begin{itemize}
            \item La faisabilité de l'AGI
            \item L'utilité de l'AGI
        \end{itemize}
    \item Discussions souvent guidées par :
        \begin{itemize}
            \item Intérêts économiques
            \item Marketing et course technologique
        \end{itemize}
\end{itemize}
\end{block}

\vspace{0.3cm}

\textbf{Rester critique et factuel face aux annonces sur l'AGI}

\end{frame}

\begin{frame}{Conclusion}

\textbf{Résumé :}

\vspace{0.3cm}

\begin{itemize}
    \item L'IA générative est une \textcolor{blue}{technologie puissante}
    \item Limitée à des \textbf{tâches spécifiques}
    \item Nombreux bénéfices mais aussi \textcolor{red}{défis éthiques et techniques}
    \item L'AGI reste un \textbf{objectif lointain}
        \begin{itemize}
            \item Sans consensus sur sa faisabilité ou son utilité
        \end{itemize}
\end{itemize}

\end{frame}

\begin{frame}{Questions ouvertes}

\begin{enumerate}
    \item Comment maximiser les bénéfices tout en minimisant les risques de l'IA générative ?

    \vspace{0.3cm}
    \item Faut-il réguler le développement de l'IA générative ? Si oui, comment ?

    \vspace{0.3cm}
    \item Que pensez-vous des prédictions sur l'AGI : objectif réaliste ou utopie/dystopie lointaine ?
\end{enumerate}

\end{frame}

